% Job posting: https://ca.linkedin.com/jobs/view/senior-computer-vision-engineer-at-musashi-ai-north-america-4462025492
\documentclass[11pt]{article}
\usepackage[left=1in,top=1in,right=1in,bottom=1in]{geometry}
\usepackage[hidelinks]{hyperref}
\usepackage{setspace}
\usepackage[T1]{fontenc}
\usepackage{newtxtext,newtxmath}
\ifdefined\pdfgentounicode
  \input{glyphtounicode}
  \pdfgentounicode=1
\fi
\setstretch{1.1}
\pagenumbering{gobble}
\begin{document}
\pagestyle{empty}

\noindent Nishal Stanislaus Silva, PhD \\
Toronto, ON \\
nishal.silva@hotmail.com \,|\, +1 613 200 2030 \\
nishal.xyz \,|\, linkedin.com/in/nishal-silva \,|\, github.com/ns2max

\vspace{1.5em}

\noindent Dear Hiring Manager,

\vspace{1em}

\noindent For 5.5 years at MAS Holdings, South Asia's largest apparel manufacturer, I built and ran computer-vision systems on live production floors, not in a lab. I designed an automated multilingual care-label QC system using OpenCV template and feature matching with an explainable defect overlay, iterating it from a handheld scanner to an industrial camera on a conveyor with PLC-controlled reject -- cutting inspection time 99.5\% while keeping a human-in-the-loop path for borderline cases. Alongside it, I built a loom-side fabric structural-defect detection system using a microscopic camera array that cut operator headcount by half. Reading the Senior Computer Vision Engineer opening at Musashi AI, this is the exact discipline the role calls for: getting machine vision out of a demo and onto a factory floor where it has to run reliably, fast, and cheaply, every shift.

\vspace{1em}

\noindent Those 5.5 years taught me what it takes to move a CV system from prototype to deployed production: camera and sensor integration, real-time inference under hard constraints, and the unglamorous work of making a system that a line operator trusts. My PhD in Information Engineering and Computer Science at the University of Trento added a layer of rigor on top -- frozen-protocol benchmarking, ablation studies, and quantified accuracy/latency/CPU trade-offs -- that I used to build and publish a real-time detector running at 14 ms on a Raspberry Pi 4 (F1 0.76, 74x faster than a DTW baseline). Across both settings I have worked as a systems engineer as much as a modeler: expert C++17 for real-time inference engines, Python for training and evaluation, and hands-on experience with high-frequency IoT ingestion, embedded ARM deployment, and CI/CD.

\vspace{1em}

\noindent I don't know Musashi AI's specific manufacturing vertical or the details of its robotics integrations from the posting alone, but the core discipline -- production-grade machine vision under real factory constraints, deployed on embedded hardware -- is precisely what I have spent my career doing, and it transfers directly.

\vspace{1em}

\noindent I am a Canadian Permanent Resident and do not require sponsorship. I am based in Toronto, open to relocating to Waterloo, and available immediately. I would welcome the chance to discuss how my industrial computer-vision background can contribute to your team.

\vspace{1.5em}

\noindent Sincerely, \\
Nishal Stanislaus Silva

\end{document}
